\begin{textsample}{POS Dim 1 – global\_north\_2024\_04 – Score 27.00 – global\_north\_2024\_04/107456944.txt}  \label{ex:f1_pos_160}
Biden officials worry that Israeli response to Iran ' s attack may trigger \textbf{wider} war President Joe Biden has privately expressed concern that Israeli Prime Minister Benjamin Netanyahu is trying to drag Washington into a broader conflict, according to three people familiar with his comments.
A senior administration official said Biden told Netanyahu in a call Saturday night that the U.S. will not participate in offensive operations against Iran and that Israel should not respond by \textbf{retaliating} against Iran.
An anti-missile system operates over Ashkelon, southern Israel, early Sunday after Iran \textbf{launched} a \textbf{drone} and \textbf{missile} attack.Amir Cohen / Reuters Within hours of the \textbf{strike}, Maj.
Gen.
Mohammad Bagheri, the chief of staff of the \textbf{Iranian} armed forces, issued a statement saying the attack had"concluded and we are not willing to continue it."An Israeli official in the prime minister ' s office told NBC News on Sunday that"Israel is going to consult with all its partners, but ultimately it ' s Israel ' s decision as to what the response will be.""Israel can't allow such @ @ @ @ @ @ @ @ @ @, be it small or large,"the official said."It ' s up to the war Cabinet to decide now."U.S. officials had expected a response from Iran since Israel ' s April 1 bombing of an \textbf{Iranian} diplomatic compound in Damascus.
They raised questions about not just the severity of the threat but also how Israel would respond given an overall military strategy that a senior administration official and a senior defense official, referring to internal conversations, have characterized in interviews as poorly considered and"frenetic"and with the potential to be"catastrophically escalatory."Israel views Iran as an existential threat and has clashed with previous U.S. administrations over how to handle the regime, which doesn't officially recognize it.
Most notably, Netanyahu openly lobbied against President Barack Obama ' s landmark 2015 \textbf{nuclear} deal with Iran, which his successor, Donald Trump, pulled out of.
Iran has been on the receiving end of a concerted U.S. and Israeli campaign to degrade its military \textbf{capabilities}, with attacks killing @ @ @ @ @ @ @ @ @ @ \textbf{region} and \textbf{nuclear} scientists at home. ' Sleepwalking ' into war?
While the White House believes the Israelis aren't looking for a \textbf{wider} conflict or a direct war with Iran, particularly given the resources they have fighting in Gaza, U.S. officials can't be certain, the senior administration official said."There ' s this urgency to act, and that ' s what happened in Damascus,"the official said, referring to the April 1 \textbf{strike} on the \textbf{consular} building, which killed two generals and five officers in the \textbf{Iranian} \textbf{Revolutionary} \textbf{Guard} Corps.
The official said it ' s the same frustration U.S. officials have with the way Israel is operating in Gaza.
Fawaz Gerges, a professor of international relations and Middle Eastern politics at the London School of Economics, criticized Biden ' s ability to translate his concerns into influence over Israel ' s military decision-making."The strategy of the Biden administration has failed miserably.
Biden is sleepwalking the U.S. into another catastrophic war in the Middle East.
His overarching @ @ @ @ @ @ @ @ @ @ neighboring countries has failed,"Gerges said."Biden has failed to influence Netanyahu ' s decisions either in Gaza or towards Iran."President Joe Biden and members of his national security team Saturday.Adam Schultz / AP Benjamin Friedman, policy director of the think tank Defense Priorities, said in a statement that"the Israeli government has courted a fight with Iran, perhaps encouraged by the prospect of U.S. help in going after Iran.""Instead of talking about ' ironclad ' support for Israel, the president should have made clear the U.S. support is limited and does not extend to all circumstances,"Friedman said."War with Iran would imperil U.S. security for no obvious pay off."Defense Priorities, based in Washington, advocates \textbf{restraint} in U.S. foreign policy.
Regardless of the relative caution U.S. officials were advocating, a U.S. military official told NBC News that U.S. forces in the \textbf{region} continue to shoot down Iranian-launched \textbf{drones} \textbf{targeting} Israel."Our forces remain postured to provide @ @ @ @ @ @ @ @ @ @ the \textbf{region},"the official said.
Israeli Defense Minister Yoav \textbf{Gallant} warned in a statement Sunday that the confrontation between Iran and Israel"is not over yet....
We must be prepared for every scenario."He called Iran"a terrorist state."As for Iran, Amal Saad, a lecturer at Cardiff University ' s School of Law and Politics who specializes in the \textbf{Iran-backed} \textbf{Lebanese} political movement and \textbf{militia} Hezbollah, warned against interpreting the limited destruction caused by Iran ' s \textbf{strikes} as reason to dismiss its intentions or willingness to enter into a conflict with Israel if provoked.
The \textbf{strikes}, which closed down schools and sent thousands of Israelis into bunkers as \textbf{sirens} sounded, \textbf{rockets} streaked overhead and \textbf{fighter} \textbf{jets} roared, was deeply disruptive in Israel, and its scale showed that"Iran ' s strategic patience was exhausted,"Saad said, referring to Iran ' s tactic of military \textbf{restraint} and deterrence.
A man holds a model \textbf{artillery} shell during a celebration in Tehran on Sunday in support of @ @ @ @ @ @ @ @ @ @ Images Iran, Saad said, was humiliated by Israel ' s bombing of its diplomatic compound, and national dignity is"one of Iran ' s foundational values."It would have been impossible for Iran"to absorb that type of humiliation,"she said.
It"would mean Iran is no longer Iran."Further \textbf{strikes} from Israel could inflame an Iran already at its limit.
That willingness was clearly underlined by Bagheri, Iran ' s military chief of staff, whose statement said that"if Israel attacks our interests we will respond with force and our next operation will be much bigger."Gerges of the London School of Economics said :"The bottom line is the Middle East appears to be descending into a \textbf{wider} \textbf{regional} war.
For now, the ball is in Israel ' s court.
Israel ' s next move will determine if the \textbf{region} descends into further military \textbf{escalation} or \textbf{de-escalation}."

% matched lemmas: artillery, capability, consular, de-escalation, drone, escalation, fighter, gallant, guard, iran-backed, iranian, jet, launch, lebanese, militia, missile, nuclear, region, regional, restraint, retaliate, revolutionary, rocket, siren, strike, target, wide
\end{textsample}
